\section{La familia SAM: de la segmentación de instancias a la segmentación semántica}

\subsection{SAM 1 y SAM 2: segmentación a nivel de instancia}

SAM (Segment Anything Model) es la familia de modelos de segmentación de propósito general presentada por Meta. Tanto SAM 1 como SAM 2 admiten la segmentación a nivel de instancia, con las siguientes formas principales de entrada:

\begin{itemize}
    \item \textbf{Point Prompt}: se indica un punto en la imagen y se segmenta el objeto en el que se ubica dicho punto (por ejemplo, un punto en el centro del cuerpo de una ballena segmenta la ballena completa)
    \item \textbf{Bounding Box Prompt}: se da un cuadro delimitador y se segmenta el objeto dentro del cuadro
    \item \textbf{Automatic Mask Generation}: sin necesidad de ningún prompt de entrada, segmenta automáticamente todas las instancias de la imagen completa
\end{itemize}

\begin{knowledgebox}{Los prompts de Point/Box de SAM ya son, en sí mismos, Visual Prompting}
En sentido amplio, los prompts de Point y Bounding Box que acepta SAM ya son, en sí mismos, una forma de Visual Prompting: se indica una posición visual concreta en la imagen original para guiar la segmentación del modelo. Este diseño permite combinar SAM con otros modelos de manera flexible.
\end{knowledgebox}

La función de segmentación automática (Automatic Segmentation) de SAM 1 y SAM 2 no requiere ninguna entrada manual: basta con recibir la imagen original para segmentar todas las instancias de la imagen completa. Esta función es un componente fundamental de métodos de Visual Prompting como Set-of-Mark.

\subsection{SAM 3: segmentación semántica susceptible de prompt}

SAM 3 representa un salto cualitativo respecto a sus predecesores: admite la segmentación a nivel semántico y es promptable (susceptible de prompt):

\begin{importantbox}{Entrada multimodal de SAM 3}
SAM 3 admite las siguientes formas de entrada:
\begin{itemize}
    \item \textbf{Entrada de texto}: se ingresa una frase nominal simple (como ``cat''), y se segmentan todos los gatos de la imagen
    \item \textbf{Entrada Point/Box}: la forma de interacción clásica heredada de SAM 1/2
    \item \textbf{Entrada de imágenes de ejemplo}: se proporcionan imágenes de ejemplos positivos y negativos, y el modelo se acerca a los positivos y se aleja de los negativos
\end{itemize}
\end{importantbox}

\begin{warningbox}{Limitación de la entrada de texto de SAM 3}
La entrada de texto de SAM 3 solo admite frases nominales (noun phrase) simples; \textbf{no admite palabras de posición} (como «el gato de la izquierda» o «la taza de hasta arriba», descripciones que incluyen relaciones espaciales). Si se necesita procesar consultas complejas que incluyan descripciones de posición, hay que recurrir a un VLM para simplificar primero la descripción compleja a una frase nominal, y después entregarla a SAM 3.
\end{warningbox}

\subsection{Combinación de agente VLM + SAM 3}

Dado que SAM 3 no admite palabras de posición, se puede construir un Agent Pipeline de VLM + SAM 3 para procesar tareas complejas de localización de objetos:

\begin{enumerate}
    \item \textbf{El usuario ingresa una consulta compleja}: por ejemplo, «el niño más a la izquierda que trae puesto un chaleco azul»
    \item \textbf{El VLM simplifica la consulta}: convierte la descripción compleja en una frase nominal simple: «child in blue vest»
    \item \textbf{Segmentación con SAM 3}: a partir del texto simplificado, se segmentan todos los niños que traen puesto un chaleco azul
    \item \textbf{Asignación de números}: se asigna un número a cada resultado de segmentación
    \item \textbf{Segunda inferencia del VLM}: a partir de la imagen segmentada y numerada, el VLM puede ubicar con precisión al niño «más a la izquierda»
\end{enumerate}

La idea central de este Pipeline es: el VLM se encarga de la comprensión y el razonamiento del lenguaje, y SAM 3 se encarga de la segmentación visual; cada uno cumple su función y, en conjunto, completan tareas complejas de localización que un solo modelo difícilmente lograría.

\subsection{Resumen de la sección}

La familia SAM evolucionó de la segmentación de instancias de SAM 1/2 a la segmentación semántica susceptible de prompt de SAM 3. SAM 3 admite entradas multimodales de texto, punto, cuadro e imágenes de ejemplo, pero el texto se limita a frases nominales simples. Mediante la combinación de agente de VLM + SAM 3, se pueden procesar tareas complejas de localización que incluyen palabras de posición.
